\section{Generation Structure as Sequential Closure}
\label{sec:generations}

\subsection{The Question of Generations}

The Standard Model contains three generations of fermions.
While this structure is experimentally established,
its origin is not explained within the conventional framework.
Within the present approach, the question is reformulated:
\begin{equation}
  \boxed{
    \text{Can the number of generations be understood from closure geometry?}
  }
\end{equation}

\subsection{Structural Basis: The Triangle}

The minimal closure unit identified in Section~\ref{sec:triangular}
is the triangle, consisting of three nodes and three edges.
Closure is not a single event, but a process involving relations along these edges.
\begin{equation}
  \boxed{
    \text{Closure develops through interactions along the edges of the triangle.}
  }
\end{equation}

\subsection{Sequential Stabilization}

Closure need not occur simultaneously across all edges.
Instead, stabilization may proceed progressively:
\begin{center}
\begin{tabular}{ll}
\toprule
Stabilized edges & Structural state \\
\midrule
1 & partially stabilized \\
2 & intermediate stabilization \\
3 & fully stabilized \\
\bottomrule
\end{tabular}
\end{center}
\begin{equation}
  \boxed{\text{Closure can be interpreted as a sequential stabilization process.}}
\end{equation}

\subsection{Interpretation as Generational Structure}

This sequential stabilization suggests a correspondence
between stabilization stage and generation:
\begin{center}
\begin{tabular}{ll}
\toprule
Stabilization stage & Interpretation \\
\midrule
1 edge stabilized & first generation (lighter) \\
2 edges stabilized & second generation (intermediate) \\
3 edges stabilized & third generation (heavier) \\
\bottomrule
\end{tabular}
\end{center}
\begin{equation}
  \boxed{
    \text{Generations can be interpreted as stages of closure stabilization.}
  }
\end{equation}
This relates generational hierarchy to structural progression
rather than independent particle families.

\subsection{Why Three}

The number of stages is determined by the number of edges
in the minimal closure unit.
Since a triangle has exactly three edges:
\begin{equation}
  \boxed{\text{The number of generational stages is bounded by three.}}
\end{equation}
No additional independent stage is available within the same minimal structure.

\subsection{Non-Simultaneity and Asymmetry}

Simultaneous stabilization of all edges would correspond to
the limiting case of a perfectly symmetric configuration.
However, from the foundational principle established in Vol.~1:
\begin{equation}
  \boxed{\text{Perfect symmetry does not produce dynamics.}}
\end{equation}
Stabilization must therefore proceed asymmetrically,
with edges locking at distinct stages.
\begin{equation}
  \boxed{
    \text{Sequential stabilization is structurally favored over
    simultaneous closure.}
  }
\end{equation}

\subsection{Relation to Mass Hierarchy}

As stabilization progresses,
closure strength increases and the closure deficit $\Delta$ decreases.
Given equation~\eqref{eq:mass}:
\begin{itemize}
  \item early-stage structures $\to$ larger $\Delta$ $\to$ lower mass,
  \item later-stage structures $\to$ smaller $\Delta$ $\to$ higher mass.
\end{itemize}
\begin{equation}
  \boxed{
    \text{Mass hierarchy can be interpreted as a consequence of
    sequential closure.}
  }
\end{equation}

\subsection{Structural Interpretation}

\begin{equation}
  \boxed{
    \text{Generations are not distinct particle types, but different
    stability stages of the same underlying structure.}
  }
\end{equation}
This contrasts with the conventional view in which generations are treated
as independent copies.
